\subsection*{把生成一句话拆成一连串预测}
模型不必一次选出完整句子。给定当前上下文，它先为下一个 token 分配概率，选择一个 token 后放回上下文，再预测下一步。token 是编码单位，可能是词、字或词的一部分，不能总按完整单词理解。

一条序列的概率是沿这条路径的条件概率乘积。第二步的分布取决于第一步已经选了什么，所以不能把每步无条件概率相乘，更不能把概率相加。取负对数后，乘积转成和，这让每个位置对训练目标的贡献可以累积。

\subsection*{局部最可能不等于整句最可能}
假设第一步选甲的概率为 $0.6$、选乙为 $0.4$；在甲之后最可能的下一步概率是 $0.5$，在乙之后则为 $0.9$。贪心先选甲，完整两步最多得 $0.3$；经过乙的路径却可达 $0.36$。因此每一步取最大值，不保证整段序列概率最大。

温度改变分数经 Softmax 后的尖锐程度。低温让高分候选占比更大，高温使分布更平缓；这会改变采样行为，却不凭空增加事实知识。概率较高描述模型的统计偏好，不等于内容一定真实。

接下来要理解的核心，是上下文如何变成这些概率。Transformer 先把 token 写成向量，再让它们交换信息，最后产生词表分数。先弄清输入、条件概率和采样的关系，后续矩阵细节才有明确任务。
